% Generated draft: requires mathematical and source review.
\documentclass[11pt]{article}
\usepackage[margin=25mm]{geometry}
\usepackage{fontspec}
\setmainfont{Latin Modern Roman}
\usepackage{amsmath,amssymb,mathtools,unicode-math}
\setmathfont{Latin Modern Math}
\usepackage{enumitem,needspace,xurl,hyperref}
\setlength{\parindent}{0pt}
\setlength{\parskip}{5pt}
\setlength{\emergencystretch}{3em}
\newcommand{\N}{\mathbb{N}}
\newcommand{\Z}{\mathbb{Z}}
\newcommand{\Q}{\mathbb{Q}}
\newcommand{\R}{\mathbb{R}}
\newcommand{\C}{\mathbb{C}}
\newcommand{\F}{\mathbb{F}}
\newcommand{\A}{\mathbb{A}}
\newcommand{\PP}{\mathbb{P}}
\newcommand{\E}{\mathbb{E}}
\newcommand{\eps}{\varepsilon}
\newcommand{\dd}{\,\mathrm{d}}
\begin{document}
\section*{UnsolvedMath ID: 20002726}
\textbf{Source problem code:} catalogue entry
\subsection{Definitions and mathematical statement}
\textbf{Recoverable source target.} The primary text is only ``Problem 6.1 above.)''. The aid identifies it as a cross-reference at the tail of Problem 6.3, not a complete independent Problem 6.1. The meaningful appended remark is V. Reiner's question: is the Fomin--Reading complex \(k\)-Cohen--Macaulay in Baclawski's sense? The missing parent problem is not reconstructed from the fragment.

Let \(\Phi\) be a finite root system of positive rank \(n\), and \(m\geq1\). Write \(\Delta^m(\Phi)\) for the Fomin--Reading generalized cluster complex, of dimension \(n-1\), whose vertices are \(m\) colored copies of the positive roots and one copy of each negative simple root, with the standard generalized compatibility relation. The precise compatibility construction is not reproduced in the source record. The source's color parameter \(k\) is renamed \(m\) to separate it from a Cohen--Macaulay index.

For a finite complex \(K\) and vertex set \(A\), let \(K\setminus A\) be its induced deletion. Over a field \(\Bbbk\), \(K\) is \(r\)-Cohen--Macaulay when every deletion of fewer than \(r\) vertices is Cohen--Macaulay over \(\Bbbk\) and has the same dimension as \(K\). Cohen--Macaulay may equivalently be specified by vanishing of reduced homology of each face link in all degrees below that link's dimension, with the usual empty-face conventions.

\textbf{Reported published answer.} The aid attributes to Athanasiadis and Tzanaki the theorem that, for every such \(\Phi,m\) and every vertex set \(A\) with \(|A|\leq m\),
\[
 \Delta^m(\Phi)\setminus A
 \quad\text{is pure shellable of dimension }n-1.
\]
Consequently \(\Delta^m(\Phi)\) is \((m+1)\)-Cohen--Macaulay over every field, a stronger answer than the literal \(m\)-Cohen--Macaulay question. The aid reports that this index is sharp, not \((m+2)\)-Cohen--Macaulay. The positivity of rank excludes an empty-root-system convention where this sharpness assertion would need separate treatment. Full primary references and the proof of sharpness are not supplied.

\textbf{Reported rank-two corollary.} For \(\Phi=I_2(a)\), \(a\geq3\), the complex is a graph with \(am+2\) vertices. The aid reports that every vertex has degree \(m+1\), and that its vertex connectivity is
\[
 \kappa(\Delta^m(I_2(a)))=m+1.
\]
In particular, deleting at most \(m\) vertices leaves it connected; for every vertex \(v\), its open neighborhood \(N(v)\), of size \(m+1\), is a minimum vertex cut. Open means that \(v\) itself is not removed. This is a rank-two consequence, not a graph-connectivity replacement for the full higher-dimensional shellability theorem.

The solved label is supported in the aid specifically for Reiner's recoverable remark and the stated corollary. It does not supply the omitted full Problem 6.3 or resolve every possible question behind the broken cross-reference.
\subsection{Short English statement}
The fragment's meaningful question is Reiner's higher Cohen--Macaulay remark. The aid reports the stronger deletion-shellability theorem: removing at most \(m\) vertices preserves the original dimension and shellability, giving sharp \((m+1)\)-Cohen--Macaulayness. Its rank-two graph has connectivity and degree exactly \(m+1\).
\subsection{Sources}
\begin{itemize}
\item Local UnsolvedMath registry, permanent ID \(20002726\), catalogue entry, \texttt{aim-probability-notes.json}, \texttt{notes[167]}, workshop \emph{Braid groups, clusters and free probability}: \url{https://aimath.org/WWN/braidgroups/braidgroups.pdf}. The source item field says 6.1; the aid identifies a tail of Problem 6.3.
\item Original remark: V. Reiner, Fomin--Reading complexes and Baclawski [5] higher Cohen--Macaulayness. Supplied research attribution: Athanasiadis and Tzanaki, deletion shellability. No complete bibliography entries, title or year are supplied.
\item Supplied rank-two corollary: attempt 1, full\_solution, status at run ``solved'', rigor ``proved'', low novelty confidence; registry solved/SOLVED-IN-LITERATURE. The machine-generated aid is not peer reviewed and requires independent verification before citation; no PDF visit is claimed.
\end{itemize}
\end{document}
